\documentclass[10pt]{article}
\usepackage[margin=1in]{geometry}
\usepackage{booktabs,array}
\usepackage[hidelinks]{hyperref}
\setlength{\parindent}{0pt}\setlength{\parskip}{4pt}
\begin{document}
{\large\textbf{S1 Text. Ideas that did not work}}\\[2pt]
HOBench manuscript. This project was run by AI agents as an experiment in agent-run research (see the main text).
The list below covers the major approaches that were tried and then abandoned, corrected or found to be
uninformative. Evidence for each item is in the repository: \texttt{docs/process/RESEARCH\_TREE.md} (decisions),
\texttt{docs/process/LESSONS.md}, \texttt{docs/reviews/}, \texttt{results/ERRATA.md}, \texttt{families/ERRATA.md} and
\texttt{results/catalogue/abandoned.md}.

\section*{1. Process}
\textbf{Running many subagents in parallel.} After the design review, the PI agent launched 13 subagents at once
(10 curators and 3 engineers). Within about an hour, they used up the account's session limit, and all of them stopped
mid-task. About 3.8 hours were lost, and two later outages cost about 7 more hours. Afterwards, the PI ran at most
four or five subagents at a time and resumed stopped subagents instead of starting new ones.

\textbf{Checking the request folder late.} The ranked list of 50 paywalled papers was posted 45 minutes into the
project. A human delivered 49 items within about 1.5 hours, but the PI checked the folder only 10 hours later.
Several curators had worked without these sources in the meantime.

\textbf{Checking data sources at curation time.} The teleost 3R ohnolog tables of Braasch et al. 2016 were planned
as a genome-scale collection. During curation, a curator found that the tables were derived from Ensembl Compara gene
trees, which the protocol excludes as truth. The collection was dropped. An earlier source audit would have saved the
planning time.

\textbf{Deriving status from free text.} The evaluation script first decided whether a family was headline-eligible
by searching its curation log for phrases such as ``not headline''. A shared status line made it misclassify one
family (E07c). The catalogue's computed column replaced the text search.

\textbf{A demonstration script in a shared directory.} A tooling agent's demonstration run overwrote the directory of
freshly downloaded proteomes. They were downloaded again.

\section*{2. Reference encoding}
\textbf{Point-valued merge levels.} The first formalism gave each reference duplication a single species-tree branch.
Reviewers showed with toy cases that a correct prediction for a teleost WGD then scored 0.33, while never merging
scored 1.0. Interval-valued references replaced point values.

\textbf{Duplication at the WGD node for yeast.} The first encoding placed the yeast WGD on the branch leading to the
post-WGD clade. The hybrid-origin model~[Marcet-Houben and Gabaldón 2015] makes this wrong, so we used an interval up
to the Saccharomycetaceae root. This trade-off made most YGOB events untestable (1,390 of 1,592), because the upper
bound lies outside the family range and many copies are single genes.

\textbf{$\beta$-lactamase classes as non-homologous decoys.} The first design used $\beta$-lactamase classes A, B, C
and D as families that must never merge. Four reviewers pointed out that classes A, C and D are homologous and that
plasmid alleles are absent from reference proteomes. Pairs of non-homologous isofunctional enzymes~[Omelchenko et
al. 2010] replaced them.

\textbf{``At most one gene per species means orthologs.''} This rule is false across WGD lineages, because reciprocal
gene loss leaves paralogs that look like orthologs. The rule was restricted to synteny-based sets.

\section*{3. Metrics}
Trivial baselines and error analyses on real data exposed five defects that the toy tests had missed (\texttt{results/ERRATA.md}):
\begin{enumerate}
\item \textbf{Untestable events.} All-singletons ``recovered'' 86\% of YGOB WGD events. Fix: score only events that
both trivial baselines can fail.
\item \textbf{Flat baseline projection.} OrthoFinder orthogroups without hierarchy were first treated as singletons
below the root, so the flat baseline equalled all-singletons. Fix: project each group to every level.
\item \textbf{Leaf-level duplications.} One big HOG ``recovered'' species-specific duplications with single-gene
copies, because these events test only whether in-paralogs are grouped. Fix: report them separately.
\item \textbf{Unresolved roots in membership F1.} The metric required families with an unresolved root to form one
group, which the reference does not assert. Fix: count only separations that the reference asserts; this made
membership F1 undefined for the prokaryote panel.
\item \textbf{Summary aggregation.} Families that shared a panel-level label were collapsed into one item, so a panel
showed $n=2$ instead of $n=5$. Fix: only true collections are collapsed.
\end{enumerate}
The lesson is to run controls on real data before scoring any method.

\section*{4. Families that were abandoned or demoted}
\begin{tabular}{p{4cm}p{11cm}}
\toprule
Candidate & Reason \\
\midrule
Teleost 3R ohnologs & source tables derived from Ensembl Compara (excluded as truth) \\
Salmonid 4R (AORe/LORe) & needs several salmonids; the \textit{Coregonus} proteome downloaded empty; dropped for budget \\
Full YGOB pillars & server refused connections; the archived file is truncated; only 2 species have track data \\
Scannell 2006 per-locus losses & delivered paper and supplements contain no per-gene identifiers \\
Protein kinome (Manning 2002) & delivered tables give identifiers for human only \\
MADS-box genes & delivered sources label identifiers for 2--3 panel species \\
Nitrogenase, GAPDH transfer, GAL-cluster HGT & fewer than 4 panel species with members \\
Myosins, plant homeobox, Sox, two adaptin families & below 80\% unique mapping (old gene models, scaffold identifiers) \\
Globins & failed blind verification (intervals wider than the source supports) \\
SwissTree reference trees & started late, stopped by an outage \\
KRAB zinc fingers, chemokines, B-family polymerases & not started (budget) \\
\bottomrule
\end{tabular}

\textbf{Motor-domain sequences.} Kinesin and myosin sources give only motor domains, which fail the 80\% coverage
rule. A domain rule ($\geq$250 residues, $\geq$90\% of the domain at $\geq$98\% identity) rescued 38 kinesins but not
enough myosins (73\% mapped).

\section*{5. Experiments}
\textbf{Testing OMAmer leakage.} We planned to run FastOMA with an OMAmer database that excludes the panel species, to
test whether FastOMA benefits from OMA-derived placements. Building such a database was not feasible within the
budget. The possible advantage remains a limitation.

\textbf{Mild robustness perturbations.} Contamination with 20 proteins touched only 8--13 units. A ten times stronger
follow-up (200 proteins) still did not change reference relations. A single changed species-tree node and removal of
two species also had no effect beyond FastOMA's run-to-run noise. Stronger species-tree errors were not tested.

\textbf{Difficulty versus agreement.} An early analysis reported that harder families show lower tool agreement
(Spearman $\rho=-0.58$, $n=20$). That analysis included non-headline families. On headline items with three methods,
$\rho=-0.41$ ($p=0.19$, $n=12$), so the claim was withdrawn.

\textbf{Confirmatory test on the first two panels.} With only the yeast and prokaryote panels, the paired test had two
items per endpoint, so the smallest possible $p$ value was 0.5. The test became meaningful only after pooling all panels.

\textbf{Quality-control covariates.} OMArk scores for completeness and contamination were computed for all proteomes but
were not used in the analysis.

\textbf{Source audit by keyword.} The audit script flagged NCBI ``LOC'' identifiers in 40 \textit{X. laevis} unit
names. Manual inspection showed that they were display labels; the evidence came from the Session et al. 2016 table.
The rule was refined.
\end{document}
